%-------------------------------------------------------------------------------
% 5. ToBig's 20th Conference — Pi-FiRi: 화재 감지 지능형 로봇 시스템
%-------------------------------------------------------------------------------
\cvsection{3. ToBig's 20th Conference : Pi-FiRi — 엣지 디바이스 기반 실시간 화재 감지 로봇}

\projsummary{라즈베리파이 로봇과 LLM Agent를 결합한 실시간 자율 순찰 및 화재 대응 시스템}
\projfigL[7.0cm]{p5.jpg}

\projpart{\faBookOpen\hspace{1.3mm}배경 및 목표}
\begin{cvitems}
  \item {\textbf{배경}\enskip 지하주차장·물류창고 등 사각지대가 많은 실내는 고정형 센서만으로 신속한 대응이 어려우며, 화재 위치나 규모를 사람이 직접 판단해야 하는 위험이 존재.}
  \item {\textbf{목표}\enskip 상용 로봇 키트에 정밀 제어·환경 센서를 통합하고, LLM Agent로 화재를 실시간 탐지·분석·보고하는 '지능형 화재 감지 시스템'을 개발.}
  \item {\textbf{개발 기간}\enskip 2025.05 -- 2025.07}
  \item {\textbf{기술 스택}\enskip \pill{Python} \pill{TensorFlow} \pill{Raspberry Pi} \pill{EfficientDet-Lite1} \pill{Flask} \pill{LangChain} \pill{Solar Pro2} \pill{OpenCV} \pill{PID Control} \pill{Streamlit}}
\end{cvitems}

\projpart{\faScrewdriverWrench\hspace{1.3mm}파이프라인}
\projfigL[7.9cm]{p5pipe.png}

\projpart{\faGear\hspace{1.3mm}시스템 및 하드웨어 구성}
\begin{cvitems}
  \item {\textbf{EfficientDet-Lite1}\enskip 엣지 구동을 위해 TFLite 모델 중 속도(49ms)와 정확도(mAP 30.55\%)의 균형점인 'Lite1'을 선정하고, Roboflow Universe의 Fire\&Smoke Dataset을 수집·증강해 파인튜닝.}
  \item {\textbf{연속 감지 필터링}\enskip 단일 프레임 노이즈로 인한 오탐을 막기 위해, 0.5 임계값 이상 탐지가 1.2초 연속 유지될 때만 화재로 확정.}
  \item {\textbf{Pi-Camera}\enskip 화재·연기를 탐지하고 실시간 영상 스트림을 서버로 전송하는 로봇의 '눈' 역할.}
  \item {\textbf{GrayScale Sensor}\enskip 바닥 라인을 실시간 인식해 지정 경로를 따라가는 라인 트레이싱으로 정해진 구역을 상시 순찰.}
  \item {\textbf{DHT11 Sensor}\enskip 카메라 영상만으로는 알 수 없는 현장 온·습도를 관제 대시보드와 실시간 Q\&A에서 함께 확인.}
  \item {\textbf{ArUco Marker (OpenCV)}\enskip 로봇의 현재 위치와 화재 발생 지점을 파악하기 위해 경로 구간별로 도입.}
\end{cvitems}

\newpage
\projpart{\faMagnifyingGlass\hspace{1.3mm}설계 과정 및 수행 내용}

\stephead{1}{PID 제어 로직을 통한 주행 안정화 (하드웨어 한계 보완)}
\begin{substeps}
  \item PiCar-4WD 키트를 단순 if-else 조건문으로 제어하면, 모터 출력 불균형 때문에 주행이 좌·우로 흔들리는 헌팅(Hunting) 현상과 경로 이탈이 발생 --- 상용 키트의 하드웨어 한계를 제어 로직으로 보완해야 함.
  \item GrayScale 센서가 읽은 바닥 라인의 오차($e$)를 수치화하고, \textbf{PID 제어}로 좌·우 모터 출력을 실시간 미세 조정해 이탈 없는 라인 트레이싱 환경을 구축.
  \begin{substeps2}
    \item \textbf{I(적분) 제어} --- 누적 오차를 반영해, 모터 출력 저하로 한쪽으로 쏠리는 주행 편향을 보정.
    \item \textbf{D(미분) 제어} --- 오차의 변화율을 감지해, 조향 시 발생하는 떨림을 억제하고 부드러운 주행을 구현.
  \end{substeps2}
\end{substeps}

\stephead{2}{분산 컴퓨팅 아키텍처 설계 (엣지 연산 한계 극복)}
\begin{substeps}
  \item 엣지 환경(라즈베리파이)에서 주행 제어와 고부하 AI 추론을 동시에 수행하면 연산 지연으로 \textbf{1 FPS}까지 떨어져, 실시간 구동 자체가 불가능.
  \item 라즈베리파이와 로컬 PC를 동일한 모바일 핫스팟(WLAN)에 연결해 하나의 서브넷을 구성하고, 역할을 분리.
  \begin{substeps2}
    \item \textbf{라즈베리파이} --- Flask API 서버를 올려 카메라 영상과 센서 데이터를 스트리밍하는 역할만 담당.
    \item \textbf{로컬 PC} --- 전송받은 프레임에 EfficientDet-Lite1 추론을 수행하고 결과를 회신.
  \end{substeps2}
  \item 이 분산 구조로 처리량을 \textbf{약 10 FPS}까지 끌어올려, 주행과 탐지가 동시에 도는 실시간성을 확보.
\end{substeps}

\stephead{3}{Agentic AI 리포팅 (Solar Pro2)}
\begin{substeps}
  \item 기존 경보 시스템은 경고음만 울릴 뿐, 화재의 발생 위치나 규모는 사람이 직접 판단해야 함 --- 탐지 결과를 사람이 읽을 수 있는 상황 보고로 바꿀 필요.
  \item 로봇이 탐지한 화재 메타데이터(감지 시간·객체 크기·예측 확률·위치)를 LangChain PromptTemplate으로 동적 주입해, Solar Pro2 기반 분석 결과를 생성.
  \begin{substeps2}
    \item \textbf{4단계 위험도 경고} --- 예측 확률과 객체 크기를 종합해 [관심--주의--경계--심각]으로 분류한 자연어 메시지를 생성.
    \item \textbf{순찰 보고서 생성} --- 순찰 시작·종료 시간, 구간별 통과 로그, 화재 이벤트를 사실에 근거해 작성. \texttt{temperature=0}과 다중 제약 시스템 프롬프트로 환각을 차단하고, 같은 입력에 같은 결과를 내는 재현성을 확보해 일관된 평가가 가능한 응답 체계를 구축.
    \item \textbf{실시간 Q\&A} --- ``현재 로봇의 위치는 어디야?''와 같은 질문에, 수집된 순찰 데이터를 근거로 실시간 답변.
  \end{substeps2}
\end{substeps}

\stephead{4}{Streamlit 관제 대시보드 구축}
\begin{substeps}
  \item 실시간 영상 스트리밍, 온·습도 센서 수치, PID 파라미터 제어, 로봇 주행 상태, LLM Agent 응답을 하나의 인터페이스에 통합 시각화.
  \item 관제자가 화면 전환 없이 ``탐지 $\to$ 판단 $\to$ 조치''를 한 곳에서 수행하도록 해 관제 효율을 극대화.
\end{substeps}
